\section*{Capítulo \ref{ch:g10:chemical-species} --- Especies químicas naturales y sintéticas}

\begin{solution}{exo:g10:chemical-species:1}
Naturales: la vainillina de una vaina, la cafeína de los granos de café.
Sintéticas: la aspirina, la vainillina de fábrica.
\end{solution}

\begin{solution}{exo:g10:chemical-species:2}
Son la misma \omterm{def:g6:pure-substances-and-mixtures:species}{especie química}, y una especie tiene las mismas propiedades
sea cual sea su origen.
\end{solution}

\begin{solution}{exo:g10:chemical-species:3}
$R_f = 2.4/6.0 = 0.40$.
\end{solution}

\begin{solution}{exo:g10:chemical-species:4}
Enfría el vapor y lo vuelve a convertir en líquido. Al entrar por el
extremo inferior, el agua llena toda la camisa (y se encuentra en último
lugar con el vapor más caliente, con lo que enfría mejor).
\end{solution}

\begin{solution}{exo:g10:chemical-species:5}
\omterm{def:g2:mixing-and-dissolving:dissolve}{Disolver} la especie mucho mejor que el agua; ser \omterm{def:g6:solutions-and-solubility:miscible}{inmiscible} con el agua;
ser lo menos peligroso posible y fácil de eliminar.
\end{solution}

\begin{solution}{exo:g10:chemical-species:6}
El etanol es \omterm{def:g6:solutions-and-solubility:miscible}{miscible} con el agua: no se forman capas y no se puede
separar nada. El ciclohexano es \omterm{def:g6:solutions-and-solubility:miscible}{inmiscible} con el agua y \omterm{def:g2:mixing-and-dissolving:dissolve}{disuelve} bien
el linalol.
\end{solution}

\begin{solution}{exo:g10:chemical-species:7}
El agua queda arriba: \qty{1}{mL} de diclorometano pesa \qty{1.33}{g},
más que \qty{1}{mL} de agua, así que se hunde.
\end{solution}

\begin{solution}{exo:g10:chemical-species:8}
Funde a una temperatura demasiado baja (aspirina: \qty{135}{\celsius}) y
a lo largo de un intervalo: la muestra es impura, o no es aspirina.
\end{solution}

\begin{solution}{exo:g10:chemical-species:9}
Linalol y etanoato de linalilo (sus dos manchas están a la altura de las
dos referencias). El etanoato de linalilo subió más alto.
\end{solution}

\begin{solution}{exo:g10:chemical-species:10}
\omterm{def:g10:chemical-species:distillation}{Hidrodestilación}; añadir ciclohexano al destilado y agitar; separar las
capas; \omterm{def:g3:separating-mixtures:evaporate}{evaporar} el ciclohexano.
\end{solution}

\begin{solution}{exo:g10:chemical-species:11}
El lápiz no se \omterm{def:g2:mixing-and-dissolving:dissolve}{disuelve} en el \omterm{def:g10:chemical-species:tlc}{eluyente}; por encima del \omterm{def:g10:chemical-species:tlc}{eluyente}, las
manchas suben por la placa en lugar de \omterm{def:g2:mixing-and-dissolving:dissolve}{disolverse} en la cubeta.
\end{solution}

\begin{solution}{exo:g10:chemical-species:12}
El ciclohexano hierve a \qty{81}{\celsius}, y el linalol, a
\qty{198}{\celsius}: un calentamiento suave elimina el ciclohexano y el
linalol se queda.
\end{solution}

\begin{solution}{exo:g10:chemical-species:13}
La muestra es probablemente esa especie (el mismo $R_f$ que la natural
en la misma placa) y probablemente pura (una sola mancha, punto de
fusión neto).
\end{solution}

\begin{solution}{exo:g10:chemical-species:14}
$0.40 \times 5.5 = \qty{2.2}{cm}$. La otra, a $0.75 \times 5.5 =
\qty{4.1}{cm}$ (al milímetro), es decir, unos \qty{1.9}{cm} por encima
de A.
\end{solution}

\begin{solution}{exo:g10:chemical-species:15}
\qty{200}{mL} pueden \omterm{def:g2:mixing-and-dissolving:dissolve}{disolver} $1.6 \times 0.200 = \qty{0.32}{g}$: con
\qty{0.20}{g} no está saturado; podría \omterm{def:g2:mixing-and-dissolving:dissolve}{disolver} \qty{0.12}{g} más. Los
\qty{0.20}{g} disueltos se perderían con el agua; el ciclohexano, que
\omterm{def:g2:mixing-and-dissolving:dissolve}{disuelve} mucho mejor el linalol, recupera la mayor parte.
\end{solution}

\begin{solution}{pb:g10:chemical-species:1}
\textbf{1.} Su vapor arrastra las especies aromáticas fuera de las
flores y hasta el refrigerante.

\textbf{2.} Enfría el vapor y lo vuelve a convertir en líquido; el agua
fría entra por el extremo inferior.

\textbf{3.} \omterm{def:g6:solutions-and-solubility:miscible}{Inmiscible}; menos denso que el agua (flota).

\textbf{4.} La capa es muy fina, y parte del aceite está disuelto en el
agua o disperso en ella: una \omterm{def:g10:chemical-species:extraction}{extracción} lo recupera.

\textbf{5.} \omterm{def:g2:mixing-and-dissolving:dissolve}{Disuelve} bien el linalol, es \omterm{def:g6:solutions-and-solubility:miscible}{inmiscible} con el agua y hierve
a una temperatura baja (\qty{81}{\celsius}), así que es fácil de
eliminar.

\textbf{6.} No: el etanol se \omterm{def:g6:pure-substances-and-mixtures:mixture}{mezcla} con el agua y no se forman capas.

\textbf{7.} Arriba (\qty{0.78}{g} por mL, menos que el agua). Con
diclorometano (\qty{1.33}{g} por mL), la capa de \omterm{def:g6:solutions-and-solubility:solute}{disolvente} quedaría
abajo.

\textbf{8.} Ninguna llama cerca (inflamable); trabajar bajo una campana
extractora y con guantes, y recoger los residuos en lugar de tirarlos
por el fregadero (tóxico para los organismos acuáticos).

\textbf{9.} El ciclohexano hierve a \qty{81}{\celsius}, y el linalol, a
\qty{198}{\celsius}: un calentamiento suave solo elimina el \omterm{def:g6:solutions-and-solubility:solute}{disolvente}.

\textbf{10.} $R_f(\text{L}) = 2.4/6.0 = 0.40$; $R_f(\text{A}) = 4.5/6.0
= 0.75$.

\textbf{11.} Linalol y etanoato de linalilo.

\textbf{12.} No: contiene al menos dos especies; es una \omterm{def:g6:pure-substances-and-mixtures:mixture}{mezcla}.

\textbf{13.} Para que todas las manchas se desarrollen en las mismas
condiciones: solo así se pueden comparar las alturas.

\textbf{14.} Su $R_f$ es igual al del linalol en las mismas condiciones:
es probablemente linalol, idéntico a la \omterm{def:g10:chemical-species:natural}{especie natural}.

\textbf{15.} $1.0 \times 0.90 = \qty{0.90}{g}$.

\textbf{16.} $0.90/100 = \qty{0.90}{\percent}$.

\textbf{17.} $100/0.009 \approx 11\,000$\,g, unos \qty{11}{kg} de
flores.

\textbf{18.} Se extrae de una planta, sin ningún \omterm{def:g5:irreversible-changes:chemical-change}{cambio químico}. Pero su
linalol es la misma especie que el linalol sintético: la misma fórmula,
las mismas propiedades.

\textbf{19.} \qty{0.90}{g} de aceite esencial a partir de \qty{100}{g}
de flores.
\end{solution}
